\documentclass[a4paper,10pt]{article}
\usepackage[utf8]{inputenc}
\usepackage{geometry}
\usepackage{graphicx}
\geometry{margin=1in}

\usepackage{xcolor}
\usepackage{hyperref}
\hypersetup{colorlinks=true,linkcolor=blue,urlcolor=blue}
\urlstyle{rm}

\usepackage{attachfile2}

\usepackage{embedfile}
\usepackage{hypgotoe}
\embedfile{references/resume-references.pdf}
\embedfile{references/42Transcript.pdf}

\usepackage{array}

%------------------------------------------------

\begin{document}

\title{Attila NÉMET}
\date{}
\maketitle


\begin{minipage}[t]{0.14\textwidth}
    \vspace{0pt}
    \includegraphics[width=\textwidth]{img/myself_ID.png}
\end{minipage}
\hfill
\begin{minipage}[t]{0.66\textwidth}
    \vspace{0pt}
    \subsection*{Coordonnées}
    \begin{tabbing}
        \hspace{2.55cm} \= \hspace{3cm} \= \kill
        \textbf{Mobile} :  \> +352-661-118922 \\
        \textbf{Email} :  \> \href{mailto:attila.nemet@gmail.com}{attila.nemet@gmail.com} \\
        \textbf{LinkedIn} :   \> \href{https://www.linkedin.com/in/attilanemet}{https://www.linkedin.com/in/attilanemet} \\
        \textbf{Adresse} :   \> 29 rue Jean-Pierre Sauvage, 2514 Luxembourg \\
    \end{tabbing}
    \end{minipage}

\subsection*{Profil}
Apprenant actuel à 42 Luxembourg, je connais de l'intérieur la pédagogie par projets, l'apprentissage par les pairs et les défis spécifiques auxquels font face les apprenants (Piscines, évaluations, travail collaboratif). Fort de plus de deux décennies d'expérience internationale chez Ericsson dans plus de 20 pays, j'ai développé de solides compétences en communication interculturelle, en accompagnement d'équipes diversifiées et en gestion de situations complexes. Titulaire d'un Master de l'Université du Luxembourg, je souhaite mettre mon parcours atypique et ma connaissance approfondie de 42 au service de l'accompagnement socio-éducatif des apprenants.

\subsection*{Formation}
\begin{itemize}
    \item \textbf{En cours} : 42 School, Luxembourg\\
    \href{gotoe:embedded=references/42Transcript.pdf}{Apprentissage par projets}, fondamentaux en C et C++, spécialisation visée en IA/LLM
    \item \textbf{Septembre 2024} : Université du Luxembourg\\
    Master en Technologies Spatiales et Business (niveau 7 CEC)
    \item \textbf{Juin 1991} : Université Technique de Cluj-Napoca, Roumanie\\
    Électronique \& Télécommunication (5 ans)
\end{itemize}

\subsection*{Compétences clés pour l'accompagnement socio-éducatif}
\begin{itemize}
    \item \textbf{Connaissance de 42 et de ses méthodes :} En tant qu'apprenant actif à 42 Luxembourg, je vis quotidiennement la pédagogie par les pairs, les Piscines et le système d'évaluation. Cette expérience directe me permet de comprendre les défis des apprenants et d'offrir un accompagnement adapté et crédible.
    \item \textbf{Accompagnement et travail d'équipe multiculturel :} Plus de 20 ans d'expérience dans plus de 20 pays
    (\href{gotoe:embedded=references/resume-references.pdf}{resume-references.pdf}),
    avec des équipes diversifiées. Habitué à écouter, accompagner et soutenir des collègues de cultures et de profils variés. Expérience de mentorat et de transmission de connaissances auprès de collaborateurs juniors.
    \item \textbf{Communication et gestion de situations complexes :} Expérience en tant que Technical Lead et Business Analyst, avec une capacité éprouvée à communiquer efficacement, à résoudre des conflits et à maintenir un climat de travail positif au sein d'équipes pluridisciplinaires.
    \item \textbf{Compréhension du marché IT et du numérique :} Connaissance approfondie du secteur des technologies de l'information, des tendances en programmation, IA/ML et des compétences recherchées sur le marché de l'emploi IT.
    \item \textbf{Insertion professionnelle :} Expérience concrète du monde professionnel IT international, permettant d'orienter et de conseiller les apprenants dans leur préparation à l'emploi.
    \item \textbf{Outils numériques :} Maîtrise des outils bureautiques et numériques (Git, VS Code, outils collaboratifs).
\end{itemize}

\subsection*{Langues}
\begin{tabbing}
    \hspace{0.55cm} \= \hspace{2.55cm} \= \hspace{3cm} \= \kill
    \> \textbf{Anglais} :  \> C2 -- Courant \\
    \> \textbf{Roumain} :  \> C2 -- Courant \\
    \> \textbf{Français} :   \> B2 -- Intermédiaire supérieur \\
    \> \textbf{Allemand} :   \> B1 -- Intermédiaire \\
    \> \textbf{Hongrois} : \> L1 -- Langue maternelle \\
\end{tabbing}

\subsection*{Compétences techniques}
\begin{tabbing}
    \hspace{0.55cm} \= \hspace{4.55cm} \= \hspace{3cm} \= \kill
    \> \textbf{Langages} :  \> C/C++, Python, Perl, LaTeX, LLM \\
    \> \textbf{Protocoles \& APIs} :  \> XML, JSON, numpy, pandas, matplotlib \\
    \> \textbf{Outils} :   \> Visual Studio Code, Jupyter Notebook, Git \\
    \> \textbf{IA} :   \> MS Copilot, Claude Code, MCP, LangChain \\
\end{tabbing}


\subsection*{Expériences professionnelles sélectionnées chez Ericsson}
\begin{itemize}
    \item \textbf{Avr 2015 -- Avr 2023, Budapest} : Business Analyst \& Testeur logiciel
    \begin{itemize}
        \item Conception et réalisation de tests de charge sur des routeurs et des nœuds télécom 4G/5G.
        \item Conseil et accompagnement pour le développement d'outils de configuration RAN.
    \end{itemize}
    \item \textbf{Mai 2013 -- Juin 2014, Dublin} : Technical Lead, déploiement 4G
    \begin{itemize}
        \item Développement de scripts, procédures et bases de données pour le programme de déploiement LTE.
        \item Encadrement et coordination d'équipes techniques multiculturelles.
    \end{itemize}
    \item \textbf{Fév 2011 -- Déc 2012, Düsseldorf} : Ingénieur développement de procédures
    \begin{itemize}
        \item Développement et exécution de procédures de modernisation RNC en réseau opérationnel, sans accès à des nœuds de test. Exécution sans incidents majeurs.
    \end{itemize}
    \item \textbf{Oct 2004 -- Jan 2010, EMEA} : Ingénieur intégration réseau senior
    \begin{itemize}
        \item Déploiement et intégration de réseaux 3G dans plusieurs pays (Égypte, Monténégro, Slovénie, Hongrie).
    \end{itemize}
    \item \textbf{Jan 1996 -- Sep 2004} : Divers postes d'ingénieur télécom
    \begin{itemize}
        \item Missions dans plus de 15 pays (Japon, Chine, Russie, Irlande, Suède, Tunisie, etc.).
        \item Adaptation constante à des environnements culturels et techniques variés.
    \end{itemize}
\end{itemize}

\subsection*{Références}
\begin{tabbing}
    \hspace{0.55cm} \= \hspace{4.55cm} \= \hspace{3cm} \= \kill
    \> \href{https://www.linkedin.com/in/jorgequerol/}{\textbf{Jorge Querol}}  \> Professeur SatCom, Université du Luxembourg\\
    \> \href{https://www.linkedin.com/in/djamilaaouada/}{\textbf{Djamila Aouada}}  \> Professeur CVIA, Université du Luxembourg\\
    \> \href{https://www.linkedin.com/in/fabrizio-pastore-41226b4/}{\textbf{Fabrizio Pastore}}  \> Professeur, directeur de mémoire\\
    \> \href{https://www.linkedin.com/in/miguel-a-olivares-mendez-8a952533/}{\textbf{Miguel Olivares-Mendez}}  \> Professeur, superviseur de projet scientifique\\
    \> \href{https://www.linkedin.com/in/christian-fisch-7745ba167/}{\textbf{Christian Fisch}}  \> Professeur d'économie et entrepreneuriat\\
\end{tabbing}


\end{document}
